\documentclass[12pt, b5paper]{article}
\usepackage{ctex}

\usepackage{geometry}
\geometry{b5paper,left=2cm,right=1cm,top=2.5cm,bottom=2.5cm}
\usepackage{chemformula} %化学公式宏包
\usepackage[version=4]{mhchem} %化学公式宏包
\usepackage{siunitx} %数字，单位宏包
\sisetup{
	separate-uncertainty = true,
	inter-unit-product = \ensuremath{{}\cdot{}}
} %单位间加小圆点
\usepackage{tikz} %Tikz绘图包
\usetikzlibrary{cd}
\usetikzlibrary{chains}
\usetikzlibrary{arrows}
\usetikzlibrary{arrows.meta} %画箭头
\usetikzlibrary{commutative-diagrams} %交换图
\usetikzlibrary{positioning,automata}
\usepackage[all]{xy}
\usepackage{booktabs} %三线表
\usepackage{multirow} %合并单元格
\usepackage{makecell} %表格内强制换行
\usepackage{array}
\usepackage{booktabs} %控制表格行高
\usepackage{colortbl}  %彩色表格需要加载的宏包
\usepackage{amsmath}
\usepackage{unicode-math}
\usepackage{fontspec} %字体
\newcommand*{\cred}{\textcolor{red}}
\setmainfont{Times New Roman}
\setmathfont{XITS Math}
\setCJKmainfont{SimSun}[AutoFakeBold, AutoFakeSlant] 

\setlength{\parindent}{0pt} %无缩进
\usepackage{setspace}
\usepackage{fancyhdr} %设置页码
\usepackage{lastpage} %获取总页数
\pagestyle{fancy} %fancyhdr宏包新增的页面风格
\renewcommand\headrulewidth{0pt} %取消页眉中的装饰分割线
\fancyhf{}
\cfoot{\footnotesize 第 \thepage\ 页(共 \pageref{LastPage}页)}%当前页 of 总页数



\begin{document}


\begin{spacing}{1.625}

\centerline{{\huge \bf{河北农业大学课程考试试卷}}}

\centerline{\bf{\large 20\underline{24}年春（\quad ）/秋（\ \surd \ ）学期}}

{\bf{考试科目：\underline{\quad 化学反应工程 \quad} \hfill 姓名：\underline{ \qquad \qquad} \hfill 学号：\underline{ \qquad \qquad \qquad \qquad}}}

{\bf{学院：\underline{\quad 理工系\quad } \hfill 专业班级：\underline{\qquad \qquad \qquad \qquad} \hfill 卷别：\cred{A} \hfill 考核方式：\cred{开卷}}}

（注：考生务必将答案写在答题纸上，标清楚大小题号，写在本试卷上无效）

\bigskip


{\bf{一、简答题（共40分）}}

1. （10分）在同样的操作条件下，完成相同的生产任务时，为什么多釜串联操作釜式反应器的生产能力大于单釜操作反应器？多釜串联是否串联级数越多越好？

\bigskip

2. （10分）绝热温升的物理意义是什么？如何计算？

\bigskip

3. （10分）液相平行反应：
\[
	\ch{A + B -> P} \qquad r_\mathrm{P} = k_1c_\mathrm{A}c_\mathrm{B}^{0.3} \ \mathrm{kmol\cdot m^{-3}\cdot min^{-1}}
\]
\[
	\ch{A + B -> Q} \qquad r_\mathrm{Q} = k_2c_\mathrm{A}^{0.5}c_\mathrm{B}^{1.3} \ \mathrm{kmol\cdot m^{-3}\cdot min^{-1}}
\]
P为目标产物，请选择一种合适的反应器型式及加料方式。

\bigskip

4. （10分）一级不可逆气固相催化连串反应：
\[
	\ce{A ->[k_1] B ->[k_2] D}
\]
请列举影响目的产物B的收率的因素，及其具体影响。


\newpage
{\bf{二、计算题(共60分)}}

1. (40分)己二酸（A）和己二醇（B）以等摩尔比缩聚反应生产醇酸树脂。反应温度70℃，催化剂为\ch{H2SO4}。实验测得动力学方程式为$r_\mathrm{A} = kc_\mathrm{A}^2$\ \si{kmol/(L.h)}。其中，速率常数$k = 118.2$ \ \si{L/(kmol.h)}，反应物的初始浓度$c_\mathrm{A0} = 0.004$\ \si{kmol/L}。每天处理2400 kg己二酸，要求最终转化率为80\%。

(1) 反应在间歇操作的釜式反应器中进行，每批辅助时间为1 h，求反应体积。

(2) 反应在全混流反应器中进行，求反应体积。

(3) 反应在活塞流反应器中进行，求反应体积。

(4) 比较不同反应器反应体积的差异，并分析原因。


\vspace{4em}
2. (20分) 某直径为6 mm的球形催化剂，其真密度为 \SI{3.60}{g/cm^3}，颗粒密度为 \SI{1.65}{g/cm^3}，比表面积为 \SI{200}{m^2/g}。使用此催化剂在0.1013 MPa、773 K下进行丁烷脱氢制备丁烯的等温气固相催化反应
\[
    \ch{C4H10 ->[Catalyst] C4H8 + H2}
\]
反应为针对丁烷的一级反应，反应速率方程为
\[
    r_\mathrm{A} = k_\mathrm{w} c_\mathrm{A} \qquad \si{mol/(g.s)}, \qquad k_\mathrm{w} = \SI{0.92}{cm^3/(g.s)}
\]
催化剂外表面对气相的传质系数$k_{\mathrm{G}} a_{\mathrm{m}} = \SI{0.23}{cm^3/(g.s)}$，进料流量为\SI{2}{m^3/min}，进料为纯丁烷。

(1) 若曲节因子为2.5，试计算内、外扩散均有影响时的总有效因子。

(2) 为消除内扩散的影响，请给出你的建议。


\end{spacing}
\end{document}
